% !PRINTED 504
\weberpage{504}
%<<CONT
once more by $F(t)$, that is, $N(t-\omega)$ is divisible by $F(t)^2$. By continuing this method of inference one recognises that $N(t-\omega)$ must be a power of $F(t)$, and that consequently the degree of $F(t)$ is a divisor of $n$ (Bd. I, \S.\,144).

We sum this up once more in the theorem:

8. \gesperrt{Every functional in $\Omega$ is the root of one and only one irreducible equation $F(t) = 0$ in $\overline{R}$ and $N(t-\omega)$ is a power of $F(t)$.}

\begin{center}
\S.\ 137.

\gesperrt{Integral Functionals.}
\end{center}

1. \textbf{Definition:} \gesperrt{A rational functional shall be called integral if its absolute value is an integer.}

From this definition it follows first of all that the sum, the difference and the product of two, and consequently also of arbitrarily many, integral rational functionals are again integral. For the product this is an immediate consequence of the theorem \S.\,136, 3. But in order to carry out the proof for the sum and the difference, we represent two integral rational functionals $A_1, A_2$ by fractional functions in such a way that they receive a primitive form as common denominator:
\[
A_1 = a_1\,\frac{E_1}{E}, \quad A_2 = a_2\,\frac{E_2}{E}.
\]
Then
\[
A_1 \pm A_2 = \frac{a_1 E_1 \pm a_2 E_2}{E},
\]
since now $a_1, a_2$ are integers, the divisor of the integral function $a_1 E_1 \pm a_2 E_2$ is the absolute value of $A_1 \pm A_2$. This is an integer, and hence $A_1 \pm A_2$ is an integral functional.

For constant rational functionals, that is, for rational numbers, Definition 1. gives the \gesperrt{integral rational numbers.}

We further set up, just as in \S.\,133 for the numbers, the following definition of the integral functionals of the field $\Omega$.

2. \textbf{Definition:} \gesperrt{A functional $\omega$ from $\Omega$ is called integral if it is the root of an equation}
\[
\Phi(t) = t^m + A_1 t^{m-1} + A_2 t^{m-2} + \cdots + A_m = 0
\]
%CONT>>
